\scanpage{39}
\begin{center}-7-\end{center}

and such that if $N$ is any non-zero submodule of $I$, then $M\cap N\neq(0)$.

Eckmann and Schopf have shown that every module has an injective hull, unique up to non-unique isomorphism [``Uber injektive Moduln'', Archiv. der. Math. \uline{4} (1953)]. We recall the following theorem, due to Matlis [``Injective modules over noetherian rings'', Pac. J. Math. \uline{8} (1958) 511-528]:

\uline{Theorem 3.1}: Let $A$ be a noetherian ring, and let $I$ be an injective $A$-module. Then $I$ can be written uniquely as a direct sum $I=\bigoplus I_\alpha$, where each $I_\alpha$ is the injective hull of the residue field $k(\mathfrak p_\alpha)$ of a suitable prime ideal $\mathfrak p_\alpha$ of $A$.

We also have information about the structure of one of these injective hulls [Matlis, ibid ]:

\uline{Proposition 3.2}: Let $A$ be a noetherian ring, let $\mathfrak p$ be a prime ideal of $A$, and let $I$ be an injective hull of the residue field $k(\mathfrak p)$ of $\mathfrak p$. Then one can write $I$ as a direct limit of submodules
\[
0\subseteq E^1\subseteq E^2\subseteq\cdots
\]
where each $E^i$ is an Artin module over the local ring $A_{\mathfrak p}$, and where, for each $i$,
\[
\dim_{k(\mathfrak p)}(E^{i+1}/E^i)=\dim_{k(\mathfrak p)}\bigl((\mathfrak p^i/\mathfrak p^{i+1})\otimes_A k(\mathfrak p)\bigr).
\]
\mspage{38}
